\documentclass[10pt, aspectratio=169]{beamer}
% Preámbulo modular para presentaciones Beamer (Modelación Estadística)
% Diseño y paleta institucional del Tecnológico de Monterrey

\documentclass[aspectratio=169,xcolor={dvipsnames,table}]{beamer}

\usepackage[utf8]{inputenc}
\usepackage[spanish,mexico]{babel}
\usepackage{amsmath,amssymb,amsfonts,mathtools}
\usepackage{graphicx}
\usepackage{listings}
\usepackage{booktabs}
\usepackage{tikz}
\usetikzlibrary{matrix,arrows,decorations.pathmorphing,shapes.geometric,positioning}

% Paleta Institucional Tecnológico de Monterrey
\definecolor{TecAzul}{HTML}{0039A6}       % Azul TEC: color primario institucional
\definecolor{TecAzulOscuro}{HTML}{1A2E51} % Azul marino oscuro
\definecolor{TecRojo}{HTML}{EC2661}       % Rojo de acento (paleta miTec)
\definecolor{TecGris}{HTML}{646464}       % Gris neutro para comentarios y subtítulos
\definecolor{TecFondoBloque}{HTML}{F4F6F9}% Fondo suave para bloques

% Configuración de Tema Beamer
\usetheme{Boadilla}
\usecolortheme[named=TecAzulOscuro]{structure}
\setbeamercolor{alerted text}{fg=TecRojo}
\setbeamercolor{palette primary}{bg=TecAzulOscuro,fg=white}
\setbeamercolor{palette secondary}{bg=TecAzul,fg=white}
\setbeamercolor{palette tertiary}{bg=TecAzulOscuro!80!black,fg=white}
\setbeamercolor{title}{fg=TecAzulOscuro,bg=TecFondoBloque}
\setbeamercolor{frametitle}{fg=TecAzulOscuro,bg=TecFondoBloque}

% Bloques personalizados elegantes
\setbeamercolor{block title}{bg=TecAzulOscuro,fg=white}
\setbeamercolor{block body}{bg=TecFondoBloque,fg=black}
\setbeamercolor{block title alerted}{bg=TecRojo,fg=white}
\setbeamercolor{block body alerted}{bg=TecRojo!8,fg=black}
\setbeamercolor{block title example}{bg=TecAzul,fg=white}
\setbeamercolor{block body example}{bg=TecAzul!8,fg=black}

% Redondear bordes y añadir sombra sutil
\setbeamertemplate{blocks}[rounded][shadow=true]
\setbeamertemplate{navigation symbols}{}

% Pie de página limpio con número de diapositiva
\setbeamertemplate{footline}{
  \leavevmode%
  \hbox{%
  \begin{beamercolorbox}[wd=.7\paperwidth,ht=2.25ex,dp=1ex,left]{author in head/foot}%
    \hspace*{2ex}\insertshortauthor~~(\insertshortinstitute) \hfill \insertshorttitle
  \end{beamercolorbox}%
  \begin{beamercolorbox}[wd=.3\paperwidth,ht=2.25ex,dp=1ex,right]{date in head/foot}%
    \insertshortdate{}\hspace*{2em}
    \insertframenumber{} / \inserttotalframenumber\hspace*{2ex}
  \end{beamercolorbox}}%
  \vskip0pt%
}

% Configuración de Listings de Python para visualización pedagógica de código
\lstset{
    language=Python,
    basicstyle=\ttfamily\footnotesize,
    keywordstyle=\color{TecAzulOscuro}\bfseries,
    stringstyle=\color{TecRojo},
    commentstyle=\color{TecGris}\itshape,
    showstringspaces=false,
    breaklines=true,
    frame=lines,
    rulecolor=\color{TecAzulOscuro!30},
    backgroundcolor=\color{TecFondoBloque},
    aboveskip=0.5em,
    belowskip=0.5em,
    literate={á}{{\'a}}1 {é}{{\'e}}1 {í}{{\'i}}1 {ó}{{\'o}}1 {ú}{{\'u}}1
             {Á}{{\'A}}1 {É}{{\'E}}1 {Í}{{\'I}}1 {Ó}{{\'O}}1 {Ú}{{\'U}}1
             {ñ}{{\~n}}1 {Ñ}{{\~N}}1 {¿}{{?`}}1 {¡}{{!`}}1
}

% Metadatos institucionales por defecto
\title[Modelación Estadística]{Modelación Estadística para Ciencia de Datos e Ingeniería}
\author[J. Castillo Colmenares]{Juliho David Castillo Colmenares}
\institute[Tec de Monterrey]{Tecnológico de Monterrey \\ Escuela de Ingeniería y Ciencias}
\date{\today}

% Comandos y macros estadísticas compartidas para las presentaciones Beamer (Metropolis)

% Conjuntos numéricos básicos
\newcommand{\R}{\mathbb{R}}
\newcommand{\N}{\mathbb{N}}
\newcommand{\Q}{\mathbb{Q}}
\newcommand{\Z}{\mathbb{Z}}
\newcommand{\set}[1]{\left\{ #1 \right\}}
\newcommand{\abs}[1]{\left|#1\right|}
\newcommand{\norm}[1]{\left\| #1 \right\|}

% Comandos estadísticos y probabilísticos del curso
\newcommand{\Var}{\operatorname{Var}}
\newcommand{\cov}{\operatorname{Cov}}
\newcommand{\corr}{\rho}
\newcommand{\comb}[2]{\begin{pmatrix} #1 \\ #2 \end{pmatrix}}
\newcommand{\s}{\sigma}
\newcommand{\particion}{\mathfrak{P}}
\newcommand{\card}[1]{\#\left( #1 \right)}
\newcommand{\seq}[3]{\set{#1}_{#2}^{#3}}
\newcommand{\E}{\mathbb{E}}
\newcommand{\Prob}{\mathbb{P}}

% Entornos pedagógicos y cajas en español académico natural
\newenvironment{discusion}{%
  \begin{alertblock}{Pregunta de Discusión}%
}{%
  \end{alertblock}%
}

\newenvironment{reflexion}{%
  \begin{alertblock}{Reflexión para el Aula}%
}{%
  \end{alertblock}%
}

\newenvironment{paradoja}{%
  \begin{alertblock}{Paradoja de Exploración}%
}{%
  \end{alertblock}%
}

\newenvironment{motivacion}{%
  \begin{exampleblock}{Motivación: Ciencia de Datos en la Práctica}%
}{%
  \end{exampleblock}%
}

\newenvironment{retoclase}{%
  \begin{block}{Reto Rápido en Equipo}%
}{%
  \end{block}%
}


\title[CDF Continua y Cuantiles]{Sección 04.02: CDF Continua y Cuantiles}
\subtitle{Función de Distribución Acumulada, Inversión y Método de Monte Carlo}
\author[J. Castillo Colmenares]{Juliho Castillo Colmenares}
\institute{Tecnológico de Monterrey}
\date{\vspace{-1.2cm}}

\begin{document}

% Slide 1: Portada
\begin{frame}[plain]
  \titlepage
\end{frame}

% Slide 2: Hoja de Ruta
\begin{frame}{Hoja de Ruta --- Capítulo 04: Variables Aleatorias Continuas}
  \footnotesize
  ¿Dónde nos encontramos en el desarrollo modular del capítulo? \pause
  \begin{itemize}\itemsep=0.08em
    \item \textbf{04.01} Función de Densidad (PDF) y Soporte Continuo \pause
    \item \textbf{04.02} Función de Distribución Acumulada (CDF) Continua y Cuantiles \emph{(Hoy)} \pause
    \item \textbf{04.03} Esperanza Matemática, Varianza y LOTUS Continuo \pause
    \item \textbf{04.04} Distribución Uniforme Continua \pause
    \item \textbf{04.05} Distribución Normal y Puntaje $Z$ \pause
    \item \textbf{04.06} Distribución Exponencial y Procesos sin Memoria \pause
    \item \textbf{04.07} Distribuciones Gamma, Beta y Weibull
  \end{itemize}
  \vspace{-0.05cm}
  \begin{block}{Objetivo de la Sesión}
    Construir la Función de Distribución Acumulada continua $F(x) = P(X \le x)$ a partir de la PDF, dominar sus propiedades axiomáticas, calcular probabilidades por diferencias de CDF, definir cuantiles como inversa $F^{-1}(p)$, y aplicar el método de inversión para generar variables aleatorias.
  \end{block}
\end{frame}

% Slide 3: Motivación
\begin{frame}{De la PDF a la CDF: Integración como Operación Inversa}
  \footnotesize
  \begin{motivacion}[¿Por qué la CDF además de la PDF?]
    La PDF codifica la densidad local de probabilidad, pero las preguntas prácticas suelen referirse a umbrales acumulados: ``¿cuál es la probabilidad de que el tiempo de espera supere los 5 minutos?'' o ``¿cuál es el percentil 95 del consumo eléctrico?'' La CDF responde estas preguntas de manera directa vía integración, y su inversa $F^{-1}$ define los \emph{cuantiles} de la distribución.
  \end{motivacion}

  \vspace{0.15cm}
  \pause
  \begin{itemize}\itemsep=0.1cm
    \item \textbf{Probabilidades acumuladas:} $F(x) = P(X \le x) = \int_{-\infty}^{x} f(t)\,dt$. \pause
    \item \textbf{Cuantiles:} $q_{p} = F^{-1}(p)$ es el valor bajo el cual ocurre la probabilidad $p$. \pause
    \item \textbf{Generación Monte Carlo:} $X = F^{-1}(U)$ con $U \sim U(0, 1)$ produce muestras exactas de $X$.
  \end{itemize}
\end{frame}

% Slide 4: Definición Formal de CDF
\begin{frame}{Definición Formal de la CDF Continua}
  \footnotesize
  La CDF de una variable aleatoria continua $X$ con PDF $f(x)$ se define como: \pause

  \vspace{0.1cm}
  \begin{block}{CDF Continua}
    $$F(x) = P(X \le x) = \int_{-\infty}^{x} f(t)\,dt, \quad x \in \mathbb{R}$$
    Por el Teorema Fundamental del Cálculo, $f(x) = F'(x)$ donde $F$ sea diferenciable. La CDF resume toda la información distribucional.
  \end{block}

  \vspace{0.1cm}
  \pause
  \begin{alertblock}{Propiedades Axiomáticas}
    \begin{itemize}\itemsep=0.05cm
      \item \textbf{Acotada:} $0 \le F(x) \le 1$ para todo $x$.
      \item \textbf{Límites:} $\lim_{x \to -\infty} F(x) = 0$ y $\lim_{x \to +\infty} F(x) = 1$.
      \item \textbf{Monótona no decreciente:} $a < b \implies F(a) \le F(b)$.
      \item \textbf{Continua por la derecha:} $\lim_{h \to 0^{+}} F(x + h) = F(x)$.
    \end{itemize}
  \end{alertblock}
\end{frame}

% Slide 5: Cálculo de Probabilidades
\begin{frame}{Cálculo de Probabilidades por Diferencias de CDF}
  \footnotesize
  Para una variable continua, las probabilidades de intervalos se calculan como diferencias de la CDF: \pause

  \vspace{0.1cm}
  \begin{block}{Probabilidad de un Intervalo}
    $$P(a \le X \le b) = F(b) - F(a) = \int_{a}^{b} f(t)\,dt$$
    Como $P(X = a) = P(X = b) = 0$ para variables continuas, los símbolos $<$ y $\le$ son intercambiables: $P(a \le X \le b) = P(a < X < b) = P(a \le X < b) = P(a < X \le b)$.
  \end{block}

  \vspace{0.1cm}
  \pause
  \begin{alertblock}{Cola Derecha y Cola Izquierda}
    \begin{itemize}\itemsep=0.05cm
      \item $P(X > a) = 1 - F(a)$ (cola derecha o complemento).
      \item $P(X \le b) = F(b)$ (cola izquierda directa).
      \item Para colas combinadas: $P(a \le X \le b) = F(b) - F(a^{-})$ donde $F(a^{-}) = \lim_{x \to a^{-}} F(x)$.
    \end{itemize}
  \end{alertblock}
\end{frame}

% Slide 6: Cuantiles
\begin{frame}{Cuantiles: Inversa de la CDF $q_{p} = F^{-1}(p)$}
  \footnotesize
  El \emph{cuantil} de orden $p$ es el valor $q_{p}$ tal que $F(q_{p}) = p$: \pause

  \vspace{0.1cm}
  \begin{block}{Función Cuantil}
    $$q_{p} = F^{-1}(p) = \inf\{x \in \mathbb{R} : F(x) \ge p\}, \quad p \in (0, 1)$$
    La función $F^{-1}$ es la inversa generalizada de $F$. Si $F$ es estrictamente monótona, la inversa clásica está bien definida.
  \end{block}

  \vspace{0.1cm}
  \pause
  \begin{alertblock}{Cuantiles de Uso Común}
    \begin{itemize}\itemsep=0.05cm
      \item \textbf{Mediana:} $m = q_{0.5}$, divide la distribución en dos mitades de probabilidad 50\%.
      \item \textbf{Percentiles:} $q_{p}$ para $p \in \{0.01, 0.05, 0.10, 0.25, 0.50, 0.75, 0.90, 0.95, 0.99\}$.
      \item \textbf{Mediana de Exponencial($\lambda$):} $m = (\ln 2)/\lambda$.
    \end{itemize}
  \end{alertblock}
\end{frame}

% Slide 7: Método de Inversión
\begin{frame}{Método de Inversión: Generación de Muestras Exactas}
  \footnotesize
  El \emph{método de inversión} permite generar muestras de cualquier distribución continua a partir de uniformes $U \sim U(0, 1)$: \pause

  \vspace{0.1cm}
  \begin{block}{Teorema de Inversión}
    Si $U \sim U(0, 1)$ y $F$ es una CDF continua y estrictamente monótona, entonces $X = F^{-1}(U)$ tiene CDF $F$.
    \begin{proof}
      $P(X \le x) = P(F^{-1}(U) \le x) = P(U \le F(x)) = F(x)$ por la monotonía de $F^{-1}$ y la uniformidad de $U$.
    \end{proof}
  \end{block}

  \vspace{0.1cm}
  \pause
  \begin{alertblock}{Ejemplos de Inversión}
    \begin{itemize}\itemsep=0.05cm
      \item Exponential($\lambda$): $X = -\ln(1-U)/\lambda \approx -\ln U/\lambda$.
      \item Uniforme($a, b$): $X = a + (b-a)U$.
      \item Cauchy: $X = \tan(\pi(U - 1/2))$.
    \end{itemize}
  \end{alertblock}
\end{frame}

% Slide 8: Prueba de Kolmogorov-Smirnov
\begin{frame}{Prueba de Kolmogorov-Smirnov: Bondad de Ajuste No Paramétrica}
  \footnotesize
  La prueba de \emph{Kolmogorov-Smirnov} compara la CDF empírica $F_{n}$ con una CDF teórica $F_{0}$ usando el supremo de la diferencia: \pause

  \vspace{0.1cm}
  \begin{block}{Estadístico KS}
    $$D_{n} = \sup_{x \in \mathbb{R}} |F_{n}(x) - F_{0}(x)|$$
    Bajo $H_{0}$ (los datos siguen $F_{0}$), $\sqrt{n} D_{n} \xrightarrow{d} K$ donde $K$ es la distribución de Kolmogorov con función de distribución:
    $$P(K \le x) = 1 - 2\sum_{k=1}^{\infty}(-1)^{k-1}e^{-2k^{2}x^{2}}, \quad x > 0.$$
  \end{block}

  \vspace{0.1cm}
  \pause
  \begin{alertblock}{Aplicaciones}
    \begin{itemize}\itemsep=0.05cm
      \item Verificación de normalidad de residuos en regresión.
      \item Validación de simulaciones de Monte Carlo vs. distribuciones teóricas.
      \item Comparación de dos muestras vía la prueba KS de dos colas.
    \end{itemize}
  \end{alertblock}
\end{frame}

% Slide 9: Aplicaciones
\begin{frame}{Aplicaciones en Inferencia y Simulación}
  \footnotesize
  La CDF continua y los cuantiles son herramientas fundamentales en análisis cuantitativo: \pause

  \vspace{0.15cm}
  \begin{itemize}\itemsep=0.12cm
    \item \textbf{Intervalos de Confianza:} $\text{IC}_{95\%} = [\hat{\theta} \pm z_{0.025} \cdot \text{SE}(\hat{\theta})]$ usa $q_{0.025} = -1.96$ y $q_{0.975} = 1.96$. \pause
    \item \textbf{Simulación Monte Carlo:} Generación de variables aleatorias vía inversión para propagación de incertidumbre. \pause
    \item \textbf{Pruebas de Hipótesis No Paramétricas:} Kolmogorov-Smirnov y Anderson-Darling para verificar ajuste de modelos. \pause
    \item \textbf{Análisis de Riesgo:} Value-at-Risk (VaR) y Conditional VaR definidos como cuantiles de la distribución de pérdidas.
  \end{itemize}
\end{frame}

% Slide 12A-1: Lab Python (Bloque 1)
\begin{frame}[fragile]{Laboratorio en Python: Validación de CDF y Cuantiles}
  \scriptsize
  Verificación de propiedades axiomáticas de la CDF y cálculo de cuantiles vía inversión numérica (\texttt{04.02\_continuous\_cdf.py}):
  \vspace{0.02cm}
  {\lstset{basicstyle=\fontsize{4.5pt}{5.5pt}\selectfont\ttfamily}
  \lstinputlisting[firstline=18, lastline=43]{../../code/04_variables_aleatorias_continuas/04.02_continuous_cdf.py}}
\end{frame}

% Slide 12A-2: Lab Python (Bloque 2)
\begin{frame}[fragile]{Laboratorio en Python: Método de Inversión}
  \scriptsize
  Implementación de $X = F^{-1}(U)$ para Exponential y Uniform con `scipy.stats.kstest`:
  \vspace{0.02cm}
  {\lstset{basicstyle=\fontsize{4pt}{5pt}\selectfont\ttfamily}
  \lstinputlisting[firstline=50, lastline=75]{../../code/04_variables_aleatorias_continuas/04.02_continuous_cdf.py}}
\end{frame}

% Slide 12B: Lab Python (Bloque 3)
\begin{frame}[fragile]{Laboratorio en Python: Test KS y Cuantiles Empíricos}
  \scriptsize
  Aplicación de Kolmogorov-Smirnov y comparación de cuantiles empíricos vs. teóricos:
  \vspace{0.02cm}
  {\lstset{basicstyle=\fontsize{4.5pt}{5.5pt}\selectfont\ttfamily}
  \lstinputlisting[firstline=85, lastline=110]{../../code/04_variables_aleatorias_continuas/04.02_continuous_cdf.py}}
\end{frame}

% Slide 12C: Salida Terminal
\begin{frame}[fragile]{Laboratorio en Python: Salida en Terminal}
  \scriptsize
  Salida estándar generada al ejecutar el script del laboratorio en Python:
  \vspace{0.02cm}
  \begin{block}{Salida Estándar en Consola}
    \fontsize{4pt}{5pt}\selectfont
    \begin{verbatim}
=== Block 1: CDF Validation & Quantile Inversion ===
CDF boundary properties:
  Uniform(0,1): F(-inf)=0 | F(+inf)=1.000000
  Exponential(2.0): F(-inf)=0 | F(+inf)=1.000000
  Normal(0,1): F(-inf)=0 | F(+inf)=1.000000
Monotonicity: all True
Quantile inversion: all match=True

=== Block 2: Inverse Transform Sampling ===
Exponential(lambda=2):
  Empirical mean: 0.4980 (theoretical: 0.5000)
  Empirical var:  0.2465 (theoretical: 0.2500)
  KS test: stat=0.001974, p-value=0.8299
Uniform(0, 5):
  Empirical mean: 2.5067 (theoretical: 2.5000)
  KS test: stat=0.002729, p-value=0.4450

=== Block 3: Quantile-Based Statistics & KS Test ===
KS statistic for n=50 Uniform(0,1) samples: D=0.1420
  Critical value (alpha=0.05): 0.1921
  SciPy KS p-value: 0.2417
  Decision: Fail to reject H0
95% CI for median: [0.4231, 0.6110]
    \end{verbatim}
  \end{block}
\end{frame}

% Slide 13A: Ejercicio Nivel 1 (Enunciado)

% Slide 13B: Ejercicio Nivel 1 (Resolución)

% Slide 14A: Ejercicio Nivel 2 (Enunciado)

% Slide 14B: Ejercicio Nivel 2 (Resolución)

% Slide 15A: Ejercicio Nivel 3 (Enunciado)

% Slide 17: Cierre
\begin{frame}{Cierre de Sección y Síntesis sobre CDF y Cuantiles}
  \begin{reflexion}[La CDF como Herramienta de Inferencia]
    La CDF continua $F(x) = P(X \le x)$ codifica toda la información distribucional y es la herramienta central para calcular probabilidades acumuladas y definir cuantiles. Su inversa $F^{-1}$ es la base de la inferencia basada en percentiles y del método de inversión para generación Monte Carlo.
  \end{reflexion}

  \vspace{0.2cm}
  \pause
  \begin{block}{Lecciones Clave de la CDF Continua}
    \begin{itemize}\itemsep=0.08cm
      \item \textbf{CDF vs. PDF:} $F$ es acumulativa; $f$ es local. $f = F'$ y $F = \int f$.
      \item \textbf{Cuantiles:} $q_{p} = F^{-1}(p)$ es la base de intervalos de confianza y pruebas no paramétricas.
      \item \textbf{Método de inversión:} $X = F^{-1}(U)$ con $U \sim U(0, 1)$ genera muestras exactas de cualquier distribución continua.
    \end{itemize}
  \end{block}
\end{frame}

% Slide 18: Perspectiva Modular
\begin{frame}{Perspectiva Modular --- El Próximo Reto}
  \scriptsize
  \begin{retoclase}[Para la próxima sesión: Esperanza, Varianza y LOTUS Continuo]
    Con la PDF y la CDF en mano, podemos ahora calcular momentos de variables continuas. La esperanza continua $\E[X] = \int x f(x)\,dx$ y la varianza $\Var(X) = \int (x - \mu)^{2} f(x)\,dx$ generalizan los conceptos discretos. El Teorema LOTUS permite calcular $\E[g(X)]$ sin derivar la distribución de $g(X)$.
  \end{retoclase}

  \vspace{0.08cm}
  \pause
  \begin{alertblock}{Preparación Recomendada}
    \begin{itemize}\itemsep=0.06cm
      \item Repasa la técnica de integración por partes y la manipulación de integrales impropias.
      \item Reflexiona sobre cómo calcular la esperanza de $g(X) = X^{2}$ o $g(X) = e^{X}$ sin conocer la PDF de $g(X)$.
    \end{itemize}
  \end{alertblock}
\end{frame}

\end{document}
